\hypertarget{chapter-10-custom-fonts-symbols-and-writing-systems}{%
\chapter{Custom Fonts, Symbols, and Writing Systems}\label{chapter-10-custom-fonts-symbols-and-writing-systems}}

The previous chapter covered using an existing OpenType font well. This chapter goes further, into territory most LaTeX books never reach: loading a font that was custom-built or commissioned rather than downloaded from a standard collection, defining new mathematical notation cleanly enough that it does not collide with anything else in a document, and, for a reader who needs it, constructing an entirely new symbol system or writing system on top of the same OpenType infrastructure the rest of this book has been using for ordinary typefaces.

\hypertarget{loading-and-subsetting-custom-fonts}{%
\section{Loading and Subsetting Custom Fonts}\label{loading-and-subsetting-custom-fonts}}

A font a designer built specifically for a project, whether commissioned professionally or produced with a font editor such as FontForge or Glyphs, loads through \texttt{fontspec} exactly the same way a standard typeface does, provided it is a valid OpenType or TrueType file: \texttt{\textbackslash{}setmainfont\{MyCustomFace\}{[}Path\ =\ ./fonts/{]}} if the font is not installed system-wide, which is the more portable choice for a project meant to be built on more than one machine, since it keeps the font file itself inside the project's own directory structure rather than depending on it being separately installed wherever the document is compiled.

Subsetting, embedding only the glyphs a document actually uses rather than a font's full glyph set, matters most for distribution rather than for compilation: a PDF with a fully embedded custom font can be large if that font contains hundreds of glyphs the document never uses, and most PDF-producing engines, LuaLaTeX included, subset embedded fonts automatically as part of PDF generation, so this is typically not something an author needs to manage by hand. Where it does become a manual concern is with fonts under a license that restricts subsetting or redistribution; a project using a commissioned or purchased font should confirm the license permits the embedding the finished PDF will contain, since font licensing terms vary considerably and are easy to overlook until a PDF is already being distributed.

\hypertarget{defining-new-mathematical-notation-cleanly}{%
\section{Defining New Mathematical Notation Cleanly}\label{defining-new-mathematical-notation-cleanly}}

Chapter 6 covered \texttt{\textbackslash{}newcommand} and macro definition generally; mathematical notation adds a specific concern beyond ordinary macros, which is collision with symbols other packages already define. \texttt{amsmath}, \texttt{mathtools}, and any specialized package a document loads all claim command names for their own symbols, and a new operator or symbol defined without checking against what is already in use risks silently overriding an existing definition, or being silently overridden by a later package load, either of which produces a document that renders differently than intended without any error being raised.

\texttt{\textbackslash{}DeclareMathOperator}, from \texttt{amsmath}, is the correct tool for a new named operator that should behave like \texttt{\textbackslash{}sin} or \texttt{\textbackslash{}lim} (upright, properly spaced from surrounding material), rather than defining it by hand with \texttt{\textbackslash{}newcommand} and \texttt{\textbackslash{}mathrm}, which technically produces similar-looking output but does not participate correctly in spacing decisions the way a properly declared operator does. For a genuinely new symbol rather than a named operator, \texttt{\textbackslash{}newcommand} is appropriate, but checking \texttt{\textbackslash{}ifdefined\textbackslash{}symbolname} before the definition, or simply grepping the project's own macro files and its package list for the intended name before committing to it, is a small habit that prevents a class of bug that is otherwise very difficult to track down, since a shadowed symbol produces no error message at all; it simply renders as whatever definition happened to win.

For a project defining a substantial body of custom notation, collecting all of it into a single, clearly named file, loaded once from the project's shared preamble as described in Chapter 3, rather than scattering symbol definitions across whichever chapter first happened to need them, keeps the full symbol vocabulary visible in one place and makes a naming collision far easier to catch before it causes a problem.

\hypertarget{building-a-private-notation-system}{%
\section{Building a Private Notation System}\label{building-a-private-notation-system}}

A research program that spans many essays and books, using a recurring, specific vocabulary of symbols and operators across all of them, benefits from treating notation as something designed once and reused deliberately, rather than invented ad hoc in whichever document happens to need a new symbol first. This is a discipline decision as much as a technical one: choosing a symbol for a recurring concept, committing to it, and using it consistently across every subsequent document is what lets a reader who has encountered the notation once recognize it immediately in a different book, rather than having to relearn a locally redefined symbol every time.

Practically, this means maintaining the shared symbol file described in the previous section not per-project but across an entire corpus, included identically at the top of every book or essay that uses the notation, with new symbols added to that shared file as the framework grows rather than defined locally and only later, if ever, reconciled with the shared vocabulary. A symbol retired or renamed in the shared file should be handled the way any breaking change is handled in a maintained codebase: with a clear record of what changed and why, since older documents referencing the previous name will otherwise fail to compile, or worse, silently compile with the wrong symbol if the old name was reused for something new rather than removed outright.

\hypertarget{custom-symbols-and-invented-writing-systems}{%
\section{Custom Symbols and Invented Writing Systems}\label{custom-symbols-and-invented-writing-systems}}

For an author who needs symbols beyond what any existing font provides, whether a small set of custom glyphs for a specific notation system or something as ambitious as a complete invented script, the underlying mechanism is the same OpenType infrastructure covered in the previous chapter, used in the opposite direction: instead of loading a font someone else built, the document loads a font built specifically to contain the needed glyphs.

Building such a font is genuinely a font-design task, done with FontForge or a comparable tool, mapping each new symbol onto a Unicode code point, either from one of Unicode's officially reserved Private Use Areas if the symbols are not meant to correspond to any standardized script, or onto an existing block if the goal is compatibility with an established but rarely supported writing system. Once built, the font loads through \texttt{fontspec} exactly as any other custom font from the previous section, and text is entered either directly as the Unicode characters the font maps its glyphs to, or, more conveniently for an author who does not want to memorize a table of private-use code points, through macros that translate a readable, ASCII-based input into the correct underlying Unicode character before typesetting, which is a straightforward application of the macro techniques from Part III to translate a document author's convenient shorthand into the code points the custom font actually expects.

This is a substantial undertaking, appropriate for a project where a distinct visual notation system is genuinely central to the work rather than a nice-to-have, but for a reader who does need it, the path from ``I need a symbol no font provides'' to ``I have typeset text in that symbol'' runs entirely through tools already introduced in this book: a font built to contain the glyph, loaded through \texttt{fontspec}, and macros to make entering it convenient.

\hypertarget{multilingual-documents}{%
\section{Multilingual Documents}\label{multilingual-documents}}

A document mixing multiple scripts, whether a Latin-script main text with occasional Greek or Cyrillic terms, or a genuinely multilingual document with substantial passages in more than one language, relies on the same Unicode-and-OpenType foundation, with a few additional considerations. \texttt{polyglossia}, the multilingual counterpart to the older \texttt{babel} package built specifically for LuaLaTeX and XeLaTeX, handles per-language hyphenation patterns, language-appropriate typographic conventions such as quotation mark style and date formatting, and switching between fonts automatically when a passage changes language, provided each language's font is specified through \texttt{\textbackslash{}setotherlanguage} alongside the main \texttt{\textbackslash{}setmainfont}.

Right-to-left and bidirectional text, needed for languages such as Arabic and Hebrew, is supported directly by both LuaLaTeX and XeLaTeX at the engine level, handling the interleaving of right-to-left and left-to-right runs within a single line correctly, which is a considerably harder typesetting problem than it appears and was one of the genuine motivations for building engines with native Unicode support in the first place, rather than relying on a system built around a Latin-script assumption from the outset. CJK typesetting, with its own line-breaking conventions distinct from space-delimited Latin text, is supported through dedicated packages (\texttt{luatexja} for LuaLaTeX, or \texttt{xeCJK} for XeLaTeX) that handle these conventions correctly once the appropriate CJK font is loaded.

A document combining several of these considerations at once, a mixed-script page with Latin body text, quoted passages in a right-to-left language, and a technical term rendered in a custom notation font from earlier in this chapter, is demanding to set up correctly, but every piece of that setup traces back to the same underlying model: Unicode code points, OpenType fonts mapped to them through \texttt{fontspec}, and an engine, LuaLaTeX or XeLaTeX, built from the ground up to handle Unicode text rather than treating it as an extension bolted onto a Latin-script core.

The next chapter returns to more conventional territory, mathematics beyond the amsmath basics, before Part V turns to managing bibliographies and reference material at the scale a book-length project actually requires.
